% Tableau des labs du cours ARGC (GES4062) dans NDAP SUKULU (propre au cours).
Le cours compte dix laboratoires, disponibles sous le code \codeUE. Les
cinq premiers portent sur l'analyse des risques, les cinq suivants sur
l'audit, les incidents, la continuité et la crise.

\begin{tableaucours}{L{2.6cm}L{6.2cm}L{1.6cm}L{1.8cm}}{%
    \tableauHeader{Lab} &
    \tableauHeader{Sujet} &
    \tableauHeader{Chapitre} &
    \tableauHeader{Opération}
}
ARGC-LAB-01 & Matrice et registre des risques & \ref{chap:matrices} & LIZBIZ \\
ARGC-LAB-02 & EBIOS RM~: menace des parties prenantes & \ref{chap:ebios} &
LIZBIZ \\
ARGC-LAB-03 & FAIR~: estimer la perte annuelle & \ref{chap:quantitatif} &
GAROUA \\
ARGC-LAB-04 & Monte-Carlo~: lire une distribution de pertes &
\ref{chap:simulation} & GAROUA \\
ARGC-LAB-05 & Indicateurs de risque et tableau de bord & \ref{chap:traiter}
& LIZBIZ \\
ARGC-LAB-06 & Échantillonnage et constats d'audit & \ref{chap:conduire} &
LIZBIZ \\
ARGC-LAB-07 & Registre des incidents~: classer et mesurer &
\ref{chap:incidents} & GAROUA \\
ARGC-LAB-08 & Analyse d'impact~: RTO, RPO, dépendances & \ref{chap:bia} &
WOURI \\
ARGC-LAB-09 & Main courante de la cellule de crise & \ref{chap:cellule} &
GAROUA \\
ARGC-LAB-10 & Communication de crise et retour d'expérience &
\ref{chap:communiquer} & MOUNGO \\
\end{tableaucours}
\vspace{-0.8em}
\begin{center}{\small\itshape Les laboratoires du cours dans NDAP SUKULU}\end{center}
